% ============================================================
% 06_evidencias.tex — Evidencias del proceso
% (Tiene mayor peso en la evaluación según las indicaciones)
% ============================================================

\section{Evidencias del Proceso}

\subsection{Estructura del repositorio}

A continuación se muestra la estructura del repositorio al final de esta semana:

\begin{lstlisting}[language=bash, caption={Estructura del repo en git.}]
so_project1/
  latex/
    preamble.tex
    refs.bib
    figures/
    avance5/        <- Este informe
      main.tex
      sections/
    sections/       <- Informe final (en construccion)
  sim/
    Makefile
    include/
      c_lock.h
      c_lock_manager.h
      core.h
    src/
      c_lock_manager.c
      core.c
      benchmark.c
\end{lstlisting}

\subsection{Diagrama de arquitectura C-Lock}

% TODO: Incluir diagrama del paper o uno propio redibujado.
% \begin{figure}[H]
%     \centering
%     \includegraphics[width=0.85\textwidth]{../../latex/figures/arquitectura_clock.pdf}
%     \caption{Arquitectura de C-Lock (adaptado del paper).}
% \end{figure}

\textit{(Diagrama pendiente — se incluirá en el siguiente avance.)}

\subsection{Extracto del código — Header principal}

% TODO: pegar aquí el header c_lock.h como evidencia de avance en código.

\begin{lstlisting}[language=C, caption={c\_lock.h — tipos y constantes del simulador.}]
/* c_lock.h - Tipos compartidos del simulador C-Lock */
#ifndef C_LOCK_H
#define C_LOCK_H

#define MAX_CORES    16
#define MAX_LOCKS    64
#define MAX_WAITERS  MAX_CORES

typedef enum { RUNNING, GATED } CoreState;

typedef struct { int owner; int waiters[MAX_WAITERS]; int nwaiters; } LockEntry;
typedef struct { int core_id; CoreState state; long cycles_active; long cycles_gated; } Core;

#endif
\end{lstlisting}

% TODO: agregar capturas de pantalla, apuntes del paper, diagramas propios, etc.
